\subsection{The full mixed six-orientation word and cycle test}

\begin{statement}{D23.C28}{Proved}{All eighteen remaining values pass the full mixed core test}
\end{statement}
Put $x=C^2$, $y=(ABA)^2$, $z=(xy)^{-1}$ and
$M=\langle x,y\rangle$ in the degree-1344 Hurwitz image.
The assignment
\[
                 \theta(x)=x,\qquad\theta(y)=z
                                                               \tag{D23.C28.1}
\]
defines an automorphism of the finite group $M$.
All eighteen raw-pentagon survivors indexed by
\[
        i\in\{1,19,23,27,53,61\},\qquad j\in\{0,1,2\}
\]
have compatible values of one commutator word on all six
peripheral orientations of this full mixed representation.
Those values satisfy the gentle two-cycle, three-cycle and
$m=-1$ three-cycle in the entire induced $C_2*C_3$ quotient.
This is a finite word-and-identity assertion.  It does not
assert that these values lift to an actual gentle element,
nor does it include a finite automorphism-generation or
pure torsion-frame conjugacy claim.

\begin{proof}
The joined marked representations $(x,y)$ and $(x,z)$ have
image order $40122452017152$, equal to $|M|$.
Both projections are onto, since $x,z$ also generate $M$.
Consequently the first projection is an isomorphism and
the second defines the automorphism (D23.C28.1).  In particular
\[
             \theta^2=\operatorname{Ad}(x),\qquad
             \kappa=\operatorname{Ad}(z)\theta
                  \text{ interchanges }y,z.               \tag{D23.C28.2}
\]
The first identity follows directly on $x,y$; the second
uses $\theta(z)=z^{-1}yz$.

For a proposed mixed value $a=f(x,y)$, put
\[
                  t=\theta(a),\quad b=ta^{-1},\quad c=t^{-1}.
                                                               \tag{D23.C28.3}
\]
Duality and the three-cycle require $b=f(y,z)$ and
$c=f(z,x)$.  Conversely these three cyclic values determine
the six oriented values once the reflection graphs are
imposed.  Indeed, as marked representations of the same
free group,
\[
 q_{xz}=\theta q_{xy},\qquad
 q_{zx}=\theta q_{yx},\qquad
 q_{zy}=\kappa q_{yz}.                                    \tag{D23.C28.4}
\]
Thus the six-kernel intersection is already the intersection
of the three cyclic kernels.  For every retained candidate
the certificate additionally checks
$\kappa(b)=b^{-1}$, so the value determined by the third
reflection graph is exactly the dual value used in the norm
calculation below.  This consistency is checked explicitly;
it is not inferred merely from membership of $(a,b,c)$.

The $M$-orbits numbered $4,8,12,13$ have total degree 492.
Restriction to their union has image order $|M|$, proving
that it is faithful.  In this smaller faithful action form
the cyclic-three image
\[
 Q=\big\langle (x,y,z),(y,z,x)\big\rangle
                      \subset M^3.
\]
Its exact order is
\[
                 |Q|=389356498710072424831060715520393216.
\]
The derived-membership test has a useful exact simplification.
Because $x,y$ have orders $6,12$ and $|M/M'|=72$, the marked
abelianization is $C_6\times C_{12}$.  A free word with
exponents $(u,v)$ has, in the first two cyclic coordinates,
the abelian values
\[
       (u\bmod6,v\bmod12),\qquad
       (-v\bmod6,u-v\bmod12).
\]
These recover both $u,v$ modulo 12.  The two generators of
$Q$ have order 12, so this proves
\[
                 Q/Q'\simeq C_{12}^2.                    \tag{D23.C28.5}
\]
In particular, a triple in $(M')^3$ belongs to $Q'$ exactly
when it belongs to $Q$.  All triples (D23.C28.3) are in
$(M')^3$.  Exact stabilizer-chain membership finds all
eighteen triples in $Q$, hence in $Q'$.  The resulting
derived order is
\[
                |Q'|=2703864574375502950215699413336064.
\]
An independent direct derived calculation and a separate
Schreier calculation over the formal source quotient
$C_{12}^2$ both gave this same order; the proof of
(D23.C28.5) avoids requiring either large calculation for
future replays.

For the full norm equations, induce the original marked
degree-1344 action from the free subgroup to $B=C_2*C_3$.
With the C22 Schreier sections, the six blocks of the
commutator-word value are
\[
       a,\quad a^{-1},\quad b,\quad t^{-1},\quad
       x^{-1}tx,\quad yb^{-1}y^{-1}.                      \tag{D23.C28.6}
\]
The last two inner corrections are verified against both
marked generators of their induced blocks.  Write $F$ for
the resulting permutation on 8064 points, and $s,r$ for
the induced torsion generators, so that
$X=(sr)^{-2}$ and $Y=sXs$.
For all $36\cdot216=7776$ values in the H7 finite fibers,
the exact equations
\[
                  (sF)^2=(F^{-1}r)^3=1
\]
hold.  The full $m=-1$ equation is checked on these same
8064 points, using the verified cyclic conjugation by
$r$ to evaluate the other orientations.  It also holds
for every value.  In particular it holds for the eighteen
remaining values considered here.  Together with the
derived-image membership and reflection consistency above,
this proves the stated finite test, with no actual
profinite lifting inference.
\end{proof}

\paragraph{Exact certificates.}
The reflection graph and its exact straight-line evaluator
are \path{agents/characters/code/hurwitz_mixed_reflection_map.py}.
This evaluates on the 29 strong generators, verifies their
reconstruction, and checks (D23.C28.2).  The faithful cyclic
image is rebuilt by
\path{agents/characters/code/hurwitz_full_mixed_s3_group.py}.
The full norm and final membership scripts are
\path{agents/characters/code/hurwitz_full_mixed_gentle_norms.py}
and \path{agents/characters/code/hurwitz_full_mixed_s3_membership.py}.
Their matching JSON/NPZ records retain every tested index
and the passing cyclic tuples.  The eighteen-element list
comes from the literal raw-pentagon record
\path{agents/pentagon/hurwitz_raw_pentagon45.npz}; its
construction and the uniqueness of the first value are
separate inputs to the main necessary reduction.
